% !TeX program = xelatex
% ============================================================
%  Содержательные файлы:
%    01-title-annotation.tex  — название работы и аннотация
%    02-goal-tasks.tex        — цель, задачи, план, риски
%    03-materials.tex         — материалы (источники) и список литературы
% ============================================================

\documentclass[a4paper,12pt,oneside]{extarticle}

% ---------- пакеты ----------
\usepackage{fontspec}
\defaultfontfeatures{Ligatures=TeX, Mapping=tex-text}
\IfFontExistsTF{Times New Roman}{\setmainfont{Times New Roman}}{%
  \IfFontExistsTF{TeX Gyre Termes}{\setmainfont{TeX Gyre Termes}}{%
    \setmainfont{Latin Modern Roman}}}
\IfFontExistsTF{Courier New}{\setmonofont{Courier New}}{\setmonofont{TeX Gyre Cursor}}

\usepackage[english,main=russian]{babel}
\usepackage{indentfirst}
\usepackage{microtype}
\usepackage[left=30mm,right=15mm,top=20mm,bottom=20mm]{geometry}
\usepackage{array,booktabs,longtable,tabularx}
\usepackage[table,dvipsnames]{xcolor}
\usepackage{enumitem}
\usepackage{amsmath,amssymb}
\usepackage[justification=centering,font=small,labelfont=bf,labelsep=period]{caption}

% ---------- оформление ----------
\setlength{\parindent}{1.25cm}
\setlength{\parskip}{0.4em plus 0.1em}
\linespread{1.3}
\clubpenalty=10000
\widowpenalty=10000
\raggedbottom

\setlist[itemize]{leftmargin=1.6cm,itemsep=1pt,topsep=2pt}
\setlist[enumerate]{leftmargin=1.9cm,itemsep=2pt,topsep=2pt}

\providecommand{\No}{\textnumero}
\newcommand{\planned}[1]{\textcolor{gray}{\textit{[план: #1]}}}
\newcommand{\done}{\textcolor{green!40!black}{\textbf{выполнено}}}
\newcommand{\inprogress}{\textcolor{orange}{\textbf{в работе}}}
\newcommand{\nextsprint}{\textcolor{blue!60!black}{\textbf{следующий спринт}}}

% ---------- заполнить один раз ----------
\newcommand{\StudentName}{Летягина Алина Алексеевна}
\newcommand{\StudentGroup}{24-Б44}
\newcommand{\SupervisorName}{Гориховский Вячеслав Игоревич}
\newcommand{\SupervisorTitle}{к.ф.-м.н., доцент кафедры системного программирования}
\newcommand{\UniversityName}{СПбГУ}
\newcommand{\FacultyName}{Математико-механический факультет}
\newcommand{\DepartmentName}{Кафедра системного программирования}
\newcommand{\CityName}{Санкт-Петербург}
\newcommand{\WorkTitle}{Разработка и программная реализация эксперимента численной оценки
бахадуровской эффективности критериев согласия в фреймворке PySATL}
\newcommand{\ReportsRepoUrl}{https://github.com/quepum/coursework-5sem}

\newcommand{\ReportPeriod}{5 семестр, спринт \No~1 (29.09.2026 -- 09.10.2026)}
\newcommand{\ReportDate}{9 октября 2026 г.}

% ============================================================
\begin{document}

% Нужна страница с оглавлением — раскомментируйте две строки ниже:
% \tableofcontents
% \clearpage

% Титульный лист
\thispagestyle{empty}
\begin{center}
    \small \MakeUppercase{\UniversityName}\\
    \small \MakeUppercase{\FacultyName}\\
    \small \MakeUppercase{\DepartmentName}
\end{center}

\vspace{1.5cm}

\begin{center}
    \large \textbf{ОТЧЁТ ПО КУРСОВОЙ РАБОТЕ}\\
    \normalsize за период: \ReportPeriod\\
    \normalsize дата: \ReportDate
\end{center}

\vfill

\begin{center}
    \Large \textbf{\WorkTitle}
\end{center}

\vfill

\noindent
\begin{tabularx}{\textwidth}{l X}
    Выполнил:             & студент группы \StudentGroup{} \StudentName \\
    Научный руководитель: & \SupervisorTitle{} \SupervisorName \\
    Работа выполняется:   & 2026/2027 учебный год (5--6 семестры) \\
    Репозиторий отчётов:  & \texttt{\ReportsRepoUrl} \\
\end{tabularx}

\vfill

\begin{center}
    \CityName, \the\year{} г.
\end{center}

\clearpage

% Аннотация
\section*{Аннотация}

Библиотека \texttt{pysatl-criterion} реализует порядка 235 критериев для 16 семейств
распределений, а \texttt{pysatl-experiment} предоставляет фреймворк Monte-Carlo
экспериментов: сегодня в нём доступны типы \texttt{critical\_value}, \texttt{power},
\texttt{time\_complexity} и \texttt{generation\_only}. При этом асимптотическая эффективность
критериев в проекте не оценивается никак: пользователь может измерить мощность на конечной
выборке и стоимость вычисления, но не может сравнить критерии по скорости убывания ошибки
первого рода при росте объёма данных. Ручной расчёт таких характеристик для десятков
критериев трудоёмок и плохо воспроизводим.

Задача работы~--- разработать численный метод оценки бахадуровской эффективности и
интегрировать его во фреймворк в виде нового типа эксперимента \texttt{bahadur\_efficiency},
переиспользующего выборки, сгенерированные из альтернатив, и уже рассчитанные предельные
распределения статистик при нулевой гипотезе. Эксперимент проходит стандартный конвейер
«генерация~--- исполнение~--- отчёт» и завершается PDF-отчётом. В границы работы не входят
вывод функций больших уклонений для статистик и реализация новых критериев; аналитические
значения эффективности предоставляются коллегой по команде и используются как эталон.

Поскольку большинство критериев библиотеки основано на статистиках с невырожденным предельным распределением,
точный наклон Бахадура для них не определён, поэтому применяется приближённый наклон. Реализация ведётся на
Python $\geq$ 3.10 с использованием NumPy, SciPy, SQLAlchemy; данные хранятся в
SQLite или PostgreSQL, вычисления распараллеливаются. Погрешность оценок количественно оценивается,
корректность подтверждается сопоставлением с аналитическими значениями.

Ожидаемые результаты: новый тип эксперимента, доведённый до состояния, когда он запускается
из CLI и даёт воспроизводимый результат; модуль численной оценки наклона и эффективности с
доверительными интервалами; матрица эффективностей для набора критериев
выбранного семейства по 3--5 альтернативам; пример конфигурации, документация и модульные тесты,
проходящие в CI.

\clearpage

% !TeX program = xelatex
% ============================================================
%  2. ЦЕЛЬ И ЗАДАЧИ РАБОТЫ
% ============================================================

\section{Цель и задачи работы}
\label{sec:goal}

\subsection{Цель}

\begin{center}
\fbox{\parbox{0.94\textwidth}{%
\textbf{Цель:} разработать численный метод оценки бахадуровской эффективности критериев
согласия, реализовать его в виде нового типа эксперимента фреймворка
\texttt{pysatl-experiment} и подтвердить корректность метода сопоставлением численных
оценок с аналитическими значениями для критериев, для которых такие значения получены.
}}
\end{center}

\subsection{Декомпозиция работы на задачи}

\begin{enumerate}[label=\textbf{\arabic*.},leftmargin=1.9cm]

  \item \textbf{Подготовительная.} \planned{5 семестр, спринт 0}
    \begin{enumerate}[label=\arabic*.\arabic*,leftmargin=1.9cm]
      \item Согласовать тему, цель и план работы с руководителем.
      \item Договориться с коллегой по команде, кто что делает: аналитические значения
            эффективности считает она, численный метод и его встраивание во фреймворк~---
            я; в каком виде она передаёт мне свои значения для сравнения.
      \item Решить с командой, в каком из двух репозиториев будет находиться код метода
            и как он попадёт к пользователям (\texttt{pysatl-criterion} подключается как
            отдельный пакет, поэтому его изменение требует публикации новой версии).
      \item Завести репозиторий отчётности, настроить автоматическую сборку PDF, написать
            аннотацию, цель и задачи.
    \end{enumerate}

  \item \textbf{Разобраться в методе и выбрать способ расчёта.} \planned{5 семестр, спринт 1--2}
    \begin{enumerate}[label=\arabic*.\arabic*,leftmargin=1.9cm]
      \item Изучить, что такое бахадуровская эффективность и как она считается;
            разобраться, почему для большинства критериев библиотеки классическая формула
            наклона не применима, и выбрать рабочий вариант~--- приближённый наклон.
      \item Определить, какие две величины нужны для расчёта и как получить каждую из них.
      \item Вывести, как из этих величин складываются наклон и эффективность~---
            относительная (критерий против критерия) и абсолютная (критерий против
            наилучшего возможного теста).
    \end{enumerate}

  \item \textbf{Доработка фреймворка под новую задачу.} \planned{5 семестр, спринт 3--4}
        Есть существующие ограничения, которые мешают реализации нового типа эксперимента; их нужно устранить.

  \item \textbf{Реализация нового типа эксперимента.} \planned{5 семестр, спринт 5--6}
    \begin{enumerate}[label=\arabic*.\arabic*,leftmargin=1.9cm]
      \item Настроить вход: новое значение типа эксперимента, его описание в конфигурации
            и проверка того, что необходимые данные уже посчитаны.
      \item Реализовать расчёт: воркер, который считает наклон и эффективность, и его
            запуск в несколько потоков.
      \item Реализовать сохранение результатов в базу и повторное использование уже
            посчитанного при перезапуске.
      \item Собрать всё вместе, чтобы эксперимент создавался и запускался из командной
            строки как существующие.
      \item Сделать отчёт: таблицу эффективностей, график зависимости от объёма выборки,
            сравнение с аналитическими значениями.
      \item Добавить пример конфигурации, документацию и тесты.
    \end{enumerate}

  \item \textbf{Проверка метода и эксперименты.} \planned{6 семестр}
    \begin{enumerate}[label=\arabic*.\arabic*,leftmargin=1.9cm]
      \item Сравнить численные оценки с аналитическими значениями коллеги; если есть
            расхождения~--- разобраться в их причинах.
      \item Проверить, что результат не меняется при увеличении числа итераций и объёма
            выборки; оценить, сколько времени занимает расчёт.
      \item Провести итоговый расчёт для выбранного семейства распределений и набора
            альтернатив, объяснить получившееся ранжирование критериев.
    \end{enumerate}

  \item \textbf{Оформление результатов и подготовка к защите.} \planned{6 семестр}
    \end{enumerate}

\end{enumerate}

\paragraph{Что не входит в задачи.} Обзор литературы отдельной задачей не является: он
ведётся непрерывно начиная со спринта~2 и оформляется файлом \texttt{03-materials.tex}.
Аналитический вывод формул и расчёт эталонных значений эффективности выполняет коллега
по команде.

% !TeX program = xelatex
% ============================================================
%  3. МАТЕРИАЛЫ (источники)
%
%  Статьи, книги, ссылки. У каждой записи четыре блока:
%    • библиографическое описание и ссылка/DOI;
%    • краткий нейросетевой обзор;
%    • зачем это нужно в работе.
%
% ============================================================

\section{Материалы}
\label{sec:materials}

\subsection{Бахадуровская эффективность}

\subsubsection*{1. Bahadur R.~R. A note on the asymptotic efficiency of statistical tests}
\textit{The Annals of Mathematical Statistics}, 1960, vol.~31, no.~2, pp.~499--503.
\textbf{Обзор.} Вводит сравнение тестов по скорости стремления достигнутого уровня
значимости к нулю; предел $-\frac{2}{n}\log L_n$ называется точным наклоном, отношение
наклонов~--- бахадуровской эффективностью.\\
\textbf{Зачем в работе.} Отправная точка подзадачи 2.1: нужно показать, что для статистик с
невырожденным предельным распределением эта конструкция вырождается, и обосновать переход к
приближённому наклону.

\subsubsection*{2. Bahadur R.~R. Stochastic comparison of tests}
\textit{The Annals of Mathematical Statistics}, 1960, vol.~31, no.~2, pp.~276--295.
DOI: 10.1214/aoms/1177705894.\\
\textbf{Обзор.} Общая схема стохастического сравнения тестов: тест эффективнее, если его
статистика быстрее уходит от критической области.\\
\textbf{Зачем в работе.} Терминологическая рамка раздела «Численный метод». Библиографические
данные подтверждены через Crossref.

\subsubsection*{3. Bahadur R.~R. Rates of convergence of estimates and test statistics}
\textit{The Annals of Mathematical Statistics}, 1967, vol.~38, no.~2, pp.~303--324.\\
\textbf{Обзор.} Связывает скорость сходимости оценок и статистик с асимптотической
эффективностью; скорость убывания хвостовых вероятностей~--- более тонкая характеристика,
чем асимптотическая дисперсия.\\
\textbf{Зачем в работе.} Аргумент во введении: почему сравнения по мощности на конечной
выборке недостаточно.

\subsubsection*{4. Nikitin Ya.~Yu. Asymptotic Efficiency of Nonparametric Tests}
Cambridge University Press, 1995. ISBN 0-521-47029-3.\\
\textbf{Обзор.} Монография: наклон Бахадура, функции больших уклонений, локальная
эффективность, эффективность непараметрических критериев согласия и независимости.\\
\textbf{Зачем в работе.} Основная теоретическая база задачи~2; источник формулировок условий
регулярности, нарушаемых при составной нулевой гипотезе.

\subsubsection*{5. Milošević B., Nikitin Ya.~Yu., Obradović M. Bahadur efficiency of EDF
based normality tests when parameters are estimated}
arXiv:2106.07437, 2021. \texttt{https://arxiv.org/abs/2106.07437}\\
\textbf{Обзор.} Для EDF-критериев нормальности при неизвестных параметрах, где
функции больших уклонений неизвестны, применяется приближённый наклон: если
$\log(1-F(t)) = -\frac{a_T t^2}{2}(1+o(1))$ и $T_n/\sqrt{n} \to b_T(\theta) > 0$, то
$c_T^{*}(\theta) = a_T b_T^2(\theta)$; эталон~--- тест отношения правдоподобий с локальным
точным наклоном $2K(\theta)$.\\
\textbf{Зачем в работе.} \textbf{Главный методологический источник.} Даёт ровно ту схему,
которую нужно реализовать численно: $a_T$ из хвоста предельного распределения (в PySATL оно
уже моделируется и сохраняется в БД), $b_T(\theta)$ из выборок по альтернативе (их уже
генерирует шаг генерации). Определяет содержание подзадач 2.1--2.3 и формат матрицы
эффективностей в подзадаче 5.3.

\subsubsection*{6. Obradović M. On asymptotic efficiency of goodness of fit tests for Pareto
distribution based on characterizations}
\textit{Filomat}, 2015, vol.~29, no.~10, pp.~2311--2324.\\
\textbf{Обзор.} Рассчитаны бахадуровские эффективности критериев согласия распределению
Парето против различных альтернатив; для каждого критерия найдены локально оптимальные
альтернативы.\\
\textbf{Зачем в работе.} Образец результата подзадачи 5.3: таблица «критерий $\times$
альтернатива $\to$ эффективность» плюс локально оптимальные альтернативы.

\subsubsection*{7. Durio A., Nikitin Ya.~Yu. Local Bahadur efficiency of some goodness-of-fit
tests under skew alternatives}
\textit{Journal of Statistical Planning and Inference}, 2003, vol.~115, no.~2,
pp.~171--179. DOI: 10.1016/S0378-3758(02)00155-6.\\
\textbf{Обзор.} Локальная бахадуровская эффективность критериев согласия при
асимметричных альтернативах.\\
\textbf{Зачем в работе.} Источник семейства асимметричных альтернатив для матрицы; важен для
подзадачи 5.2, где локальный режим наиболее неустойчив численно.

\subsubsection*{8. Akritas M.~G., Kourouklis S. Local Bahadur efficiency of score tests}
\textit{Journal of Statistical Planning and Inference}, 1988, vol.~19, no.~1,
pp.~187--199. DOI: 10.1016/0378-3758(88)90072-9.\\
\textbf{Обзор.} Локальный наклон скор-тестов; соотношение локальной и нелокальной
эффективности.\\
\textbf{Зачем в работе.} Обоснование того, что локальная и точная эффективности могут
совпадать,~--- важно для интерпретации расхождений с аналитическими эталонами
(подзадача 5.1).

\subsubsection*{9. He X., Shao Q.-M. Bahadur efficiency and robustness of studentized score
tests}
\textit{Annals of the Institute of Statistical Mathematics}, 1996, vol.~48, no.~2,
pp.~295--314. DOI: 10.1007/BF00054792.\\
\textbf{Обзор.} Выведены точные наклоны студентизированных скор-тестов; отмечена прямая
связь бахадуровской эффективности с достигнутым уровнем значимости.\\
\textbf{Зачем в работе.} Пример полного вывода наклона, на фоне которого видно, какие шаги
заменяются численной оценкой; подтверждает корректность оценки через достигнутые уровни.

\subsubsection*{10. Nikitin Ya.~Yu. Bahadur asymptotic efficiency of integral tests for
symmetry}
\textit{Journal of Soviet Mathematics}, 1984, vol.~27. DOI: 10.1007/BF01843558.\\
\textbf{Обзор.} Наклон для интегральных критериев симметрии выражается через ядро
статистики.\\
\textbf{Зачем в работе.} Часть критериев в \texttt{pysatl-criterion} относится к тому же
классу; полезно при выборе семейства для целевого прогона (задача 5.3).




\end{document}
